\documentclass[10pt,twocolumn]{article}
\usepackage[scaled]{helvet}
\renewcommand\familydefault{\sfdefault}
\usepackage[T1]{fontenc}
\usepackage[margin=0.7in]{geometry}
\usepackage{titlesec}
\usepackage{enumitem}
\usepackage{parskip}
\setlength{\columnsep}{0.24in}
\setlist[itemize]{leftmargin=1.2em, topsep=0.2em, itemsep=0.18em}
\titleformat{\section}{\large\bfseries}{\thesection}{1em}{}

\begin{document}
\title{\vspace{-1.6cm}AWS SageMaker Training}
\author{AWS ML Study Notes}
\date{}
\maketitle
\small

\section*{Overview}
SageMaker Training runs managed training jobs on ephemeral infrastructure. You provide code, data locations, and resource configuration, and SageMaker provisions the instances, streams logs, stores artifacts, and tears the instances down when the job completes.

This is the core managed primitive for repeatable training because it separates model code from the underlying host lifecycle and makes experiments easier to automate.

\section*{Core Concepts}
\begin{itemize}
\item A training job is defined by the training image, entry point script, input channels, output path, instance count, instance type, and optional distributed training settings.
\item Data can be delivered through multiple modes such as file, fast file, or pipe depending on access patterns and startup needs.
\item Managed spot training, checkpoints, profiler tooling, and metrics extraction can reduce cost and improve visibility for long or expensive jobs.
\end{itemize}

\section*{How It Works}
\begin{itemize}
\item The training script is packaged with framework specific estimators or custom containers and launched with a role that can access the required data and artifact locations.
\item SageMaker downloads or streams the data, executes the container, captures logs to CloudWatch, and writes the model artifact and any generated metrics back to S3.
\item A completed job can feed later stages such as model evaluation, registration, batch transform, or endpoint deployment.
\end{itemize}

\section*{AI and ML Relevance}
\begin{itemize}
\item Managed training is well suited for scheduled retraining, hyperparameter tuning, and experiments that need clean, isolated, reproducible execution environments.
\item Distributed training support matters when model size or dataset size outgrows a single node, especially for deep learning workloads on GPU fleets.
\end{itemize}

\section*{Operational Notes}
\begin{itemize}
\item Checkpointing is essential when using spot instances or any training run that would be expensive to restart from zero.
\item Track both infrastructure metrics and model metrics. Good GPU utilization with bad validation loss still means the experiment failed.
\item Keep containers lean and training inputs explicit so runs are reproducible and startup times do not become an invisible source of waste.
\end{itemize}

\section*{What To Remember}
\begin{itemize}
\item Training jobs are ephemeral by design. Persistent state belongs in S3 or attached managed storage, not only inside the container filesystem.
\item Input channels define what the job sees; artifact output defines what the rest of the platform can reuse.
\item If you cannot reproduce a training run from code, data version, parameters, and image version, the platform is not mature enough.
\end{itemize}

\end{document}
